\documentclass[11pt]{article}

% Packages
\usepackage{fontspec}
\setmainfont{Latin Modern Roman}
\usepackage{geometry}
\usepackage{microtype}
\usepackage{setspace}
\usepackage{amsmath,amssymb,amsthm}
\usepackage{csquotes}
\usepackage{hyperref}
\usepackage{booktabs}
\usepackage{tikz-cd}

\geometry{margin=1in}
\setstretch{1.15}

% Theorem environments
\newtheorem{definition}{Definition}
\newtheorem{proposition}{Proposition}
\newtheorem{theorem}{Theorem}
\newtheorem{remark}{Remark}

% Macros
\newcommand{\Hist}{\mathcal{H}}
\newcommand{\Future}{\mathcal{F}}
\newcommand{\Act}{\mathcal{A}}
\newcommand{\Tr}{\mathrm{Tr}}

% Title
\title{Throwing the Game:\\
Refusal, Event-Driven Cognition, and the Survival of Value\\
Beyond Autoregressive Intelligence}
\author{Flyxion}
\date{December 20, 2025\\Revised: \today}

\begin{document}

\maketitle

\begin{abstract}
Recent advances in autoregressive artificial intelligence suggest that cognition
may be reducible to the extension of statistical regularities over sequences.
On this view, intelligence consists in the compression and continuation of
history, with no need for explicit world models, commitments, or irreversibility.
This paper argues that such accounts fail to capture a structurally distinct
class of cognitive acts: refusals.

A refusal is not a preference, constraint, or optimization tradeoff. It is an
event that eliminates admissible futures and binds an agent to a history.
Where autoregressive systems preserve optionality, refusal spends it. To make
this precise, the paper introduces Spherepop, a minimal event calculus in which
agents operate over branching future spaces via irreversible operators such as
$\mathbf{Refuse}$, $\mathbf{Pop}$, $\mathbf{Bind}$, and $\mathbf{Collapse}$.

Within this framework, we prove a formal impossibility result: refusal cannot be
represented as utility maximization over a fixed triple $(\Act, U, \pi)$ without
encoding event-history as a primitive state variable. This requires explicitly
distinguishing the representing utility $U$ from the agent's recognized
instrumental evaluation functional $V$---a distinction that standard decision
theory collapses. We further prove that $\mathbf{Collapse}$ is the unique
admissibility-expanding operator in the calculus, and show that its idempotence
($\mathcal{C} \circ \mathcal{C} = \mathcal{C}$) connects Spherepop to the
theory of idempotent monads and Karoubi envelopes. Revocation is shown not to
be the inverse of refusal but to require a $\mathbf{Collapse}$ operation that
destroys historical differentiation, and ``fake refusal'' is characterized as
behavior that preserves a collapse path restoring eliminated futures.

An information-theoretic analysis sharpens a key distinction: refusal is not
Bayesian updating (conditioning on observation) but an endogenous modification
of the sample space itself---a deliberate reduction of future entropy whose cost
is conserved as constraint rather than dissipated as error. We identify three
equivalent characterizations of the same underlying phenomenon---entropy
reduction, future-space pruning, and accumulation of non-commutativity---and
show that worldhood is the obstruction to commutativity, the structure that
cannot be rearranged. Active inference is reinterpreted: refusal is incompatible
with free-energy minimization unless the generative model includes commitment
variables as latent structure. We introduce a stabilization-based halting
criterion grounded in the exhaustion of world-relevant distinctions and relate
the event algebra to IIT, assembly theory, and the phenomenology of commitment.

The central claim is that intelligence is not exhausted by prediction. It also
consists in the capacity to close worlds. Refusal is one of the events by which
such closure becomes possible.
\end{abstract}

\newpage
% ============================================================
\section{From Sequence to Event}
\label{sec:bridge}
% ============================================================

The unifying theme of this paper is a shift from sequence-based to event-based
cognition. Autoregressive systems model the extension of histories; Spherepop
models the transformation of futures. The distinction is not merely technical.
It marks the difference between systems that continue and systems that commit.

What follows should be read as an attempt to articulate this distinction at
multiple levels: formal, informational, social, and philosophical. The formal
core appears in the non-representability result of Section~\ref{sec:formal} and
the algebraic structure of Section~\ref{sec:algebra}; the remaining sections
extend these results into information theory, competition theory, active
inference, and the philosophy of mind. A synthesizing observation---that entropy
reduction, future-space pruning, and non-commutativity accumulation are
equivalent characterizations of a single phenomenon called closure---is
developed across these sections and consolidated in the worldhood analysis of
Section~\ref{sec:worldhood_unified}.

% ============================================================
\section{Introduction}
% ============================================================

Recent advances in generative artificial intelligence have challenged many
assumptions that structured twentieth-century cognitive theory. Large
autoregressive models---whether applied to language, images, audio, or
video---now produce behavior that appears coherent, context-sensitive, and
purposive without relying on explicit symbolic representations of world
structure, causal relations, or articulated goals. These systems generate
fluent dialogue, physically plausible motion, and socially legible action by
learning to extend statistical regularities in data.

The success of such systems has inspired a strong thesis: cognition may
fundamentally consist in deep autoregressive generation over learned sequences.
On this view, intelligence is not grounded in explicit internal world models but
in the ability to capture long-range dependencies that encode latent structure.
Memory exists to compensate for partial observability, and cognitive power scales
with the depth and richness of temporal context an agent can maintain. In the
limit, the argument goes, the ability to compress and extend sequential
regularities is sufficient for much of what we ordinarily treat as world
understanding \cite{barenholtz2025sequences}.

Yet this picture leaves a striking class of human behaviors unexplained. Agents
frequently and deliberately choose actions that knowingly degrade performance,
restrict future options, or incur irreversible loss. Whistleblowers sacrifice
careers to preserve integrity; artists reject lucrative contracts to protect
aesthetic autonomy; individuals refuse opportunities that would increase power at
the expense of relationships. These actions are not mistakes, nor products of
ignorance, nor mere cases of bounded rationality. They are refusals: deliberate
acts of irreversible commitment undertaken with full awareness of their cost.

This paper argues that refusal constitutes a fundamental boundary for
sequence-based accounts of cognition. Refusal is not a preference, constraint, or
auxiliary objective layered onto optimization. It is an \emph{event}: an
irreversible operation that eliminates admissible futures and binds the agent to
a specific history. While autoregressive systems preserve optionality, refusal
spends it. While sequence learning generates trajectories through possibility
space, refusal generates worlds in which certain paths are no longer available.

Methodologically, the argument is hybrid. It combines conceptual analysis
(especially in philosophy of action and commitment) with formal modeling in an
event calculus. The event calculus is Spherepop: a small set of irreversible
operators that act directly on branching future spaces. Formal results are used
primarily as constraints on interpretation. In particular, we prove a
non-representability theorem: refusal cannot be represented as ordinary utility
maximization on a fixed triple $(\Act, U, \pi)$ without augmenting the state
space with event-history. A crucial component of this proof is the distinction
between the \emph{representing} utility $U: \mathcal{O} \to \mathbb{R}$ and
the \emph{instrumental evaluation functional} $V$ the agent continues to
recognize---a distinction that standard decision theory collapses, and whose
separation is required to make the impossibility result precise. This formal
result underpins the later claims about ``strategy stealing'' and the
competitive disadvantage of value-sustaining agents \cite{carlsmith2025goodness}.

The argument proceeds in stages. First, we introduce an illustrative case in
which refusal is unusually clean: deliberate underperformance in the service of
relationship (\S\ref{sec:daryl}). Second, we characterize the autoregressive
baseline and identify its structural bias toward optionality preservation
(\S\ref{sec:ar}). Third, we develop Spherepop as a calculus with syntax,
operational semantics, and worked examples (\S\ref{sec:spherepop}). Fourth, we
distinguish refusal from preference change and related phenomena
(\S\ref{sec:prefref}), then prove the non-representability result
(\S\ref{sec:formal}). We then develop information-theoretic, competitive,
worldhood, collapse, and halting analyses, followed by a treatment of the
algebraic structure of the event calculus and its relationship to adjacent
frameworks including integrated information theory, assembly theory, and active
inference. We close with objections and open problems.

% ============================================================
\section{An Illustrative Case: Throwing the Game}
\label{sec:daryl}
% ============================================================

The 1985 film \emph{D.A.R.Y.L.} provides a particularly clear illustration of
refusal as an event rather than a preference. In the film, a child with
superhuman perceptual and motor abilities participates in a local baseball game.
He performs flawlessly. His pitch tracking, timing, and execution exceed human
capacity, resulting not merely in victory but in domination. The game ceases to
be competitive; its social meaning begins to collapse under the weight of his
competence.

This excellence produces an unexpected consequence. The child's adoptive mother,
who had previously occupied a meaningful role through care, instruction, and
support, becomes functionally redundant. Her contribution is erased not by
malice but by perfection. The local social world reorganizes itself around the
child's capability, and in doing so it removes a form of dependence that had
been constitutive of the relationship. The mother is not merely outperformed;
she is made unnecessary.

After reflection, the child deliberately begins to play poorly. His errors are
not sporadic or accidental; they are sustained and intentional. When questioned,
he explains that under certain conditions, error is more efficient than maximum
performance---specifically when the aim is relating with others. Crucially, the
explanation does not invoke ignorance, impulse, or confusion. It appeals to a
structural judgment about what continued optimization would do to the relational
field.

This is the key diagnostic point. If refusal were merely multi-objective
optimization---``win'' plus ``mother's happiness''---then it would be
representable as a standard tradeoff in a utility function. But the scene's
intelligibility rests on something stronger: the child's perfect competence
remains available and instrumentally valuable, yet is actively excluded. The
child does not discover that winning was never valuable; rather, he treats
certain victories as no longer admissible, not merely less preferred. A
trajectory---continued flawless play---is removed from the live future.

This removal is not local. It is not merely the selection of a different move at
one time-step. It changes the social topology of what can happen next: the
mother's role reappears, the game regains uncertainty, and the distribution of
attention changes. The act functions as an event in the strict sense: it closes
a branch of the future and thereby restructures what can follow.

\subsection{Counterfactual analysis}

The counterfactual ``continue playing perfectly'' is not simply ``win more.'' It
is a different world. It produces a drift toward competitive collapse (the game
becomes an exhibition), relational flattening (others become spectators), and
role elimination (care becomes redundant). If one takes seriously the thought
that relationships are sustained partly by mutual incompleteness, then perfect
performance is not merely overachievement; it is a structural solvent. The
refusal therefore has a natural interpretation as a world-preserving constraint:
it is an operation that prevents a social phase transition.

We will exploit this example as a running case: in each major section we ask
what the framework predicts about the admissible futures, the cost of eliminating
them, and the public legibility of the resulting commitment.

% ============================================================
\section{The Autoregressive Baseline}
\label{sec:ar}
% ============================================================

In order to argue that refusal is a boundary case, the baseline must be stated
precisely. We use ``autoregressive cognition'' in a broad but formal sense.

\begin{definition}[Autoregressive generator]
An autoregressive generator is a system that models a distribution over
sequences by factorization into one-step conditionals:
\[
P(x_{1:T}) = \prod_{t=1}^T P(x_t \mid x_{<t}),
\]
where $x_t$ may represent tokens, frames, actions, or observations, and where
the model's primary objective is to minimize expected predictive loss under this
factorization.
\end{definition}

This definition covers language models trained by next-token prediction,
sequence models for video and audio trained by next-step prediction or denoising
objectives, and policy models trained to imitate action sequences. The details
differ, but a shared structural commitment remains: the system is optimized to
continue producing locally coherent extensions of partial histories.

Two properties are especially relevant. First, autoregressive systems are
naturally \emph{optionality-preserving}. Their loss is defined over
continuations; ``silence'' or ``closure'' is not generally a privileged outcome
unless added explicitly. Second, they are naturally \emph{revision-friendly}.
When contradiction appears, the easiest response is to generate a continuation
that smooths it away. This is not a moral defect; it is a consequence of
training and representation.

\subsection{What autoregressive systems can do}

It is important to concede real power. An autoregressive system can model
long-range dependencies, learn implicit regularities, and generate sequences that
simulate deliberation. It can even generate \emph{performances of refusal}:
utterances of the form ``I refuse,'' ``I will not do that,'' and so on. It can
also implement refusal-like behavior if refusal is coded as a state-dependent
policy constraint supplied externally.

What it cannot do, absent additional machinery, is generate refusal as a
\emph{history-binding event}. In the D.A.R.Y.L.\ case, refusal is not merely
an action choice; it is the elimination of a class of admissible futures:
``perfect play as a live option.'' Autoregressive models can imitate such
behavior locally, but they lack an intrinsic notion of irreversible branch
elimination. Unless the training regime supplies explicit penalties for later
using the excluded action, the model's default is to treat the action as forever
re-accessible.

\subsection{Attempted refusal and the re-access problem}

Suppose an autoregressive policy $\pi_\theta(a_t \mid h_t)$ produces a
suboptimal action $a_t$ (a missed catch) in order to maintain a relationship.
If later, in a nearby context, perfect play is again locally beneficial,
$\pi_\theta$ has no intrinsic reason to treat it as inadmissible. The model can
``change its mind'' without paying a structural price. The consequence is that
refusal, when present in such systems, is typically a \emph{pattern}, not a
\emph{commitment}. This difference is exactly what Spherepop is designed to
formalize.

% ============================================================
\section{Spherepop: Syntax and Operational Semantics}
\label{sec:spherepop}
% ============================================================

To describe refusal formally, we require a non-state-based framework. The
Spherepop calculus treats agency as operating over a branching space of possible
histories. Its primitives are not states but irreversible operations that
transform future spaces.

\subsection{Future spaces}

Fix a time $t$ and an agent with a current history $h_t \in \Hist$. Let
$\Future(h_t)$ denote the set of admissible future continuations from $h_t$.
We do not assume $\Future(h_t)$ is explicitly enumerated; it is a semantic
object: the space of continuations the agent takes to be live.

\begin{definition}[Spherepop event]
A Spherepop event is an operator $E$ that maps a future space to a restricted
future space:
\[
E:\Future(h_t)\to \Future'(h_t) \subseteq \Future(h_t).
\]
An event is \emph{proper} if $\Future'(h_t)\subsetneq \Future(h_t)$.
\end{definition}

Intuitively, a proper event is characterized by exclusion: it removes at least
one admissible continuation. This exclusion is not mere preference; it is
structural. The qualification ``proper'' is necessary to distinguish genuine
events from $\mathbf{skip}$, which acts as the identity.

\subsection{Formal syntax}

We present a minimal syntax sufficient for composition.

\begin{definition}[Spherepop expressions]
Let $\mathsf{Exp}$ be the set of expressions generated by:
\[
e ::= \mathbf{skip} \mid \mathbf{Pop}(\varphi) \mid \mathbf{Bind}(\alpha \prec \beta)
\mid \mathbf{Refuse}(\alpha) \mid \mathbf{Collapse}(\psi) \mid e_1 ; e_2
\]
where $\alpha,\beta$ range over event labels or action-types, and $\varphi,\psi$
range over predicates on futures (filters) that can be evaluated against a
continuation.
\end{definition}

The sequencing operator $;$ composes events in time. $\mathbf{skip}$ is the
identity event. This syntax is intentionally modest: it does not presuppose a
full type theory or a richly structured logic of predicates. The core claim of
this paper does not depend on such elaborations. What matters is that
expressions denote operators on future spaces and that composition corresponds
to operator composition.

\subsection{Operational semantics}

We give reduction rules in terms of how an expression transforms a future
space. Write $\llbracket e \rrbracket(\Future)$ for the future space obtained
by applying event $e$.

\begin{remark}
Operationally, one can treat a future space $\Future$ as an implicit tree of
continuations. Events prune or rewrite the tree.
\end{remark}

\begin{definition}[Semantics of primitives]
Let $\Future$ be a future space. The semantics of each primitive is as follows.
$\llbracket \mathbf{skip} \rrbracket(\Future) = \Future$, leaving the future
space unchanged. Sequential composition is defined by
$\llbracket e_1 ; e_2 \rrbracket(\Future) =
\llbracket e_2 \rrbracket(\llbracket e_1 \rrbracket(\Future))$, applying events
in order. The pop operator removes futures matching a predicate:
$\llbracket \mathbf{Pop}(\varphi) \rrbracket(\Future) =
\{f\in \Future : \neg \varphi(f)\}$. Refusal removes all futures initiated by
a given action-type:
$\llbracket \mathbf{Refuse}(\alpha) \rrbracket(\Future) =
\{f\in \Future : f \text{ does not begin with action-type }\alpha\}$. The bind
operator enforces an ordering constraint:
$\llbracket \mathbf{Bind}(\alpha \prec \beta) \rrbracket(\Future) =
\{f\in \Future : \text{if } \beta \text{ occurs in } f \text{ then } \alpha
\text{ occurs earlier}\}$. Finally, Collapse quotients the future space by an
equivalence relation $\sim_\psi$ induced by predicate $\psi$, yielding a reduced
representation in which some historical differentiation is erased:
$\llbracket \mathbf{Collapse}(\psi) \rrbracket(\Future)
= \Future / \sim_\psi$.
\end{definition}

Two notes matter for later sections. First, $\mathbf{Refuse}(\alpha)$ is not
the selection of an alternative; it is deletion of a class of immediate
continuations. Second, $\mathbf{Collapse}$ is the dual operation: it sacrifices
historical differentiation to regain flexibility in a coarser space. This
duality will matter when we discuss simulated refusal versus genuine refusal:
simulated refusal behaves like a local choice; genuine refusal behaves like
deletion.

\subsection{Collapse as the unique admissibility-expanding operator}

The formal relationship between $\mathbf{Collapse}$ and all other primitives
can be stated as a sharp structural result.

\begin{proposition}[Monotone restriction outside Collapse]
\label{prop:monotone}
For every $e \in \mathsf{Evt} \setminus \{\mathbf{Collapse}\}$ and every
future space $\Future$,
\[
\llbracket e \rrbracket(\Future) \subseteq \Future.
\]
That is, every Spherepop operator except $\mathbf{Collapse}$ is
admissibility-decreasing or admissibility-preserving. $\mathbf{Collapse}$ is
the unique operator that can increase admissibility by identifying previously
distinct histories.
\end{proposition}

\begin{proof}
By inspection of the semantics. $\mathbf{skip}$ returns $\Future$ unchanged.
$\mathbf{Pop}(\varphi)$, $\mathbf{Refuse}(\alpha)$, and
$\mathbf{Bind}(\alpha \prec \beta)$ each define a subset of $\Future$ by
imposing additional conditions. Sequential composition $e_1 ; e_2$ applies two
such restrictions in succession; by induction, the result is a subset of
$\Future$. Only $\mathbf{Collapse}(\psi)$, by forming the quotient
$\Future / \sim_\psi$, can map histories into equivalence classes that reunite
previously distinguished continuations, thereby expanding access in the quotient
topology.
\end{proof}

This proposition has an immediate corollary that will be used repeatedly below:
revocation cannot be accomplished by any sequence of non-Collapse events, since
no such sequence can restore eliminated futures. Collapse is therefore not merely
an operator but a structural necessity for any account of revocation or recovery
of admissibility.

\subsection{Idempotence of Collapse}

\begin{proposition}[Idempotence]
\label{prop:idempotent}
For any Collapse equivalence relation $\sim_C$ on $\mathcal{H}$, the induced
Collapse functor $\mathcal{C}$ satisfies:
\[
\mathcal{C} \circ \mathcal{C} = \mathcal{C}.
\]
\end{proposition}

\begin{proof}
Collapse replaces each history $h$ with its equivalence class $[h]_C$. Applying
Collapse a second time to $\mathcal{H}_C = \mathcal{H}/\sim_C$ requires an
equivalence relation on equivalence classes. Since $\sim_C$ identifies histories
by future-space identity, any two $\sim_C$-classes $[h]_C$ and $[h']_C$ that
are themselves $\sim_C$-equivalent must already be equal as classes
(by transitivity of the original $\sim_C$). Therefore the second application of
$\mathcal{C}$ produces the same quotient: $\mathcal{C}(\mathcal{C}(\mathcal{H}))
= \mathcal{C}(\mathcal{H})$.
\end{proof}

\begin{remark}
Idempotent endofunctors of the form $\mathcal{C} \circ \mathcal{C} = \mathcal{C}$
are precisely the reflections in category theory, and the corresponding algebras
form a full subcategory closed under the functor. In the setting of monads, an
idempotent monad corresponds to a localization or a closure operator. The
Karoubi envelope construction, which formally adjoins splitting idempotents to a
category, provides the canonical completion. Spherepop Collapse is thus
interpretable as a Karoubi-type projection: it identifies the ``up to refusal''
equivalence classes of histories and embeds the resulting structure into a
smaller, less differentiated category. This connection makes precise the sense in
which Collapse is not arbitrary erasure but a principled localization of the
event algebra.
\end{remark}

\subsection{Worked example: formalizing the baseball case}

Let $\alpha$ be the action-type \emph{Max} (maximal performance), and let
$\beta$ be \emph{Rel} (relational maintenance, i.e.\ maintaining the mother's
meaningful role). Assume the pre-refusal future space $\Future_0$ contains
continuations in which the next play is of type \emph{Max} as well as
continuations with deliberate underperformance \emph{Err}. The key point is
that \emph{Max} is live.

We encode the refusal as the event:
\[
e_{\text{D}} := \mathbf{Refuse}(\emph{Max}).
\]
Applying it yields:
\[
\Future_1 = \llbracket e_{\text{D}} \rrbracket(\Future_0)
= \{f \in \Future_0 : f \text{ does not begin with } \emph{Max}\}.
\]

Now add a binding constraint that makes the refusal publicly and temporally
stable. For example, the agent may bind \emph{Rel} to occur before any return
to \emph{Max}:
\[
e_{\text{B}} := \mathbf{Bind}(\emph{Rel} \prec \emph{Max}).
\]
Then
\[
\Future_2 = \llbracket e_{\text{B}} \rrbracket(\Future_1)
= \{f \in \Future_1 : \emph{Max}\text{ occurs only after }\emph{Rel}\}.
\]
The combined Spherepop program is:
\[
e := e_{\text{D}} ; e_{\text{B}}.
\]
Interpretively: first exclude maximal performance in the relevant local context,
then impose an ordering constraint that makes any future re-entry into maximal
play conditional on the restoration of relational structure. This representation
lets us state precisely what the refusal does: it does not merely choose an
error; it deletes a live branch. It also shows how one can model the
``public'' aspect of refusal: a bind makes the structure legible and stable
across time.

By Proposition~\ref{prop:monotone}, the only operator that could restore
\emph{Max}-initiating futures after this program is $\mathbf{Collapse}$: that is,
a structural amnesia event that erases the distinction that made the refusal
meaningful. Revocation without Collapse is formally impossible.

\subsection{Comparison with standard frameworks}

Spherepop is not intended to replace MDPs or extensive-form games in their own
domains. It is intended to represent a phenomenon those frameworks treat only
derivatively: irreversible elimination of admissible futures as a first-class
operation.

\begin{table}[h]
\centering
\begin{tabular}{@{}llll@{}}
\toprule
Framework & Sequences & Refusal (as deletion) & Worldhood (as accumulated closure) \\
\midrule
MDP & \checkmark & \(\times\) & \(\times\) \\
POMDP & \checkmark & \(\times\) & \(\times\) \\
Extensive-form game & \checkmark & \(\circ\) & \(\circ\) \\
Autoregressive generator & \checkmark & \(\times\) & \(\times\) \\
Spherepop & \checkmark & \checkmark & \checkmark \\
\bottomrule
\end{tabular}
\caption{Expressive distinctions (schematic). Extensive-form games can represent
commitment devices by constraining strategy sets, but refusal as an intrinsic,
history-binding deletion operator is not primitive. Spherepop makes deletion
primitive, establishes Collapse as the unique expansion operator, and supports
accumulation into worldhood.}
\end{table}

The table is intentionally schematic. ``$\circ$'' indicates partial
representability via modeling tricks (e.g.\ precommitment as a move that changes
available strategies). The claim is not that other frameworks cannot simulate
refusal; it is that they do not treat refusal as primitive and therefore
naturally collapse it into preference, chance, or exogenous constraint.

% ============================================================
\section{Refusal vs.\ Preference Change}
\label{sec:prefref}
% ============================================================

A central objection is that refusal is simply multi-objective optimization or
preference update: the agent discovers that relational value outweighs winning,
and then chooses the action that maximizes this revised utility. If that were
right, Spherepop would be dispensable.

The objection gains plausibility because ordinary decision theory can represent
many superficially refusal-like patterns. If an agent values both winning and
social harmony, and learns that maximal performance harms harmony, then choosing
a suboptimal play can be utility-maximizing.

However, the refusal phenomenon at issue is structurally different. First,
refusal is temporally asymmetric. Preference change revises the ranking of
options; it does not, by itself, close options. A preference change can be
revised again at no structural cost. Refusal, by contrast, is characterized by
\emph{persistence of exclusion}: after the refusal, the excluded action remains
instrumentally valuable and remains, in many contexts, physically possible, yet
is treated as inadmissible. That persistence is what makes refusal publicly
legible as reliability.

Second, refusal is identity-binding. A preference change is a change in what
the agent wants. A refusal is a change in what the agent will allow itself to
do. This ``will not'' is not reducible to a momentary ranking; it functions as
an invariant imposed on future deliberation.

Third, refusal is essentially social in its epistemology. Others can recognize
a refusal and coordinate around it because it is not merely a local choice but
a structural alteration that persists. By contrast, preference changes are often
opaque and reversible; they do not automatically generate trust.

A useful formal criterion follows. Let $\Future_0$ be a future space and let
$\alpha$ be an action-type that is physically available and instrumentally
valuable relative to some evaluation functional $V$ (e.g.\ winning). A refusal
of $\alpha$ is an event $E$ such that $\alpha$-initiating futures are removed
from $\Future_0$ even though $V$ continues to rank them highly. In symbols:
\[
\exists f \in \Future_0 \text{ beginning with }\alpha \text{ such that } V(f)
\text{ is maximal, but } f \notin \llbracket E \rrbracket(\Future_0).
\]
This captures the distinctive feature: exclusion of a valuable future without
treating it as valueless.

In the D.A.R.Y.L.\ case, the counterfactual ``continue perfect play'' remains
valuable with respect to baseball success. The act is intelligible precisely
because the agent excludes it anyway. If the agent merely discovered that
winning was not valuable, the act would not have the same relational meaning.
It would be indifference, not refusal.

% ============================================================
\section{Formal Consequences: The Non-Representability Theorem}
\label{sec:formal}
% ============================================================

We now state and prove the non-representability result with full precision,
making explicit the distinction between the representing utility $U$ and the
agent's recognized instrumental evaluation functional $V$.

\subsection{The representation target class}

Standard decision theory represents rational choice as maximization of expected
utility induced by a triple $(\Act, U, \pi)$, where $\Act$ is a fixed,
time-indexed action set, $U: \mathcal{O} \to \mathbb{R}$ is a utility function
over outcomes in $\mathcal{O}$, and $\pi: \Hist \to \Delta(\Act)$ is a policy
mapping histories to distributions over actions. The policy $\pi$ is then
derived as:
\[
\pi(h_t) \in \arg\max_{a\in \Act} \;\mathbb{E}[U(o_{t:\infty}) \mid h_t, a].
\]
We say that a triple $(\Act, U, \pi)$ of this form \emph{represents refusal} if
it can distinguish deliberate exclusion of an action-type that remains
instrumentally high-value from ordinary preference change in which the
action-type becomes low-value.

\subsection{Separating $U$ from $V$}

A critical preliminary is the distinction between two utilities that standard
decision theory conflates. Let $V: \mathcal{F}(h_t) \to \mathbb{R}$ be the
agent's recognized \emph{instrumental evaluation functional}---the measure of
value the agent continues to consciously attribute to futures (e.g., $V(f) =
\text{win probability along }f$). Let $U: \mathcal{O} \to \mathbb{R}$ be the
\emph{representing utility} posited by the decision-theoretic model. Refusal
is the condition under which $U$ and $V$ come apart:
\[
\text{Refusal: }\quad \pi(h_t)(\alpha) = 0 \quad\text{while}\quad
\alpha \in \arg\max_{a\in \Act} \mathbb{E}[V \mid h_t, a].
\]
The agent refuses $\alpha$ even though it recognizes $\alpha$ as the best
action with respect to $V$. Standard decision theory collapses $U$ and $V$ by
assuming that whatever the policy selects is revealed as the utility-maximizing
action, eliminating the conceptual space in which refusal lives. The present
framework insists they are separate and that the gap between them is precisely
the phenomenon requiring explanation.

\subsection{Proposition and proof}

\begin{proposition}[Non-representability without history]
\label{prop:nonrep}
Let $(\Act, U, \pi)$ be a representation triple with fixed $\Act$ (not
history-dependent) and $U: \mathcal{O} \to \mathbb{R}$ defined without
reference to refusal events. Suppose an agent exhibits refusal of $\alpha$ at
time $t$: the policy assigns $\pi(h_t)(\alpha) = 0$ while $\alpha \in
\arg\max_{a \in \Act} \mathbb{E}[V \mid h_t, a]$ for the agent's recognized
instrumental functional $V$. Then no such triple can represent this refusal
without encoding event-history as a primitive state variable.
\end{proposition}

\begin{proof}[Proof sketch]
Assume for contradiction that a triple $(\Act, U, \pi)$ represents the refusal.
Since $\pi(h_t)(\alpha) = 0$, it must be that $\alpha$ is not $U$-optimal:
\[
\mathbb{E}[U \mid h_t, \alpha] < \mathbb{E}[U \mid h_t, a^*]
\]
for some $a^* \neq \alpha$. But refusal, as characterized, requires that
$\alpha$ remain $V$-optimal: the agent continues to recognize $\alpha$'s
instrumental superiority with respect to $V$. Therefore, if $U$ makes $\alpha$
strictly suboptimal, $U$ must assign disvalue to $\alpha$'s outcomes in a way
that diverges from $V$.

There are exactly two ways this divergence can arise. Either $U$ encodes a
different set of values than $V$ (in which case the agent's preferences have
changed---refusal collapses into preference change, violating the assumption),
or $U$ is conditioned on a history variable $r_t \in \{0,1\}$ recording that
a refusal event occurred (so $U = U(\cdot; r_t)$, augmenting the state from
$h_t$ to $(h_t, r_t)$). In the first case the distinction between $U$ and $V$
is erased and refusal cannot be separated from preference change. In the second
case, the action set and state description must be augmented with the commitment
variable $r_t$, which is exactly event-history. In either case, the fixed triple
$(\Act, U, \pi)$ without history augmentation is insufficient. Contradiction.
\end{proof}

\subsection{Concrete counterexample}

Return to the baseball case. Let $\alpha = \emph{Max}$, and let $V$ measure
win probability. Assume $\emph{Max}$ uniquely maximizes $V$ given the agent's
ability. The refusal consists in assigning zero probability to $\emph{Max}$ in
the relevant context while continuing to recognize that it maximizes $V$.

Any attempt to represent this in a fixed triple $(\Act, U, \pi)$ must either
(i) assign disvalue to winning when achieved via $\emph{Max}$ under $U$, thereby
redefining preferences so that $U$ and $V$ no longer diverge (collapsing refusal
into preference change), or (ii) condition $U$ on the refusal event (importing
event-history into the state). This counterexample is not about baseball; it is
about the structural role of inadmissibility. The gap between $U$ and $V$---the
representing utility and the recognized instrumental value---is the formal locus
of refusal. Closing that gap, by any means, eliminates the phenomenon.

% ============================================================
\section{Information-Theoretic Analysis}
\label{sec:info}
% ============================================================

Spherepop makes refusal a pruning operation on a future space. This naturally
admits an information-theoretic reading: refusal reduces entropy over future
possibilities. However, the conceptual precision required here is finer than
simple entropy reduction. Refusal must be distinguished from ordinary Bayesian
updating, a distinction that is central to the paper's claims about commitment
and constraint.

\subsection{Entropy over futures}

Let $\Future(h_t)$ be a set (or $\sigma$-algebra) of admissible continuations.
Let $P_t$ be the agent's subjective distribution over $\Future(h_t)$, induced
by its generative model, policy, or both. Define the entropy of the future
space as:
\[
H_t := - \sum_{f\in \Future(h_t)} P_t(f)\log P_t(f),
\]
or the analogous integral in continuous cases.

A Spherepop event $E$ produces a restricted space
$\Future'(h_t) \subseteq \Future(h_t)$. The updated distribution is the
renormalized restriction:
\[
P_t'(f) = \frac{P_t(f)\mathbf{1}_{f\in \Future'(h_t)}}{Z},
\quad
Z = \sum_{f\in \Future'(h_t)} P_t(f).
\]
The entropy after the event is $H_t' = H(P_t')$.

\subsection{Refusal is not Bayesian updating}

A subtle but crucial distinction must be made explicit before proceeding. There
are two ways in which a distribution over futures can be restricted, and they
are not equivalent.

The first is \emph{Bayesian updating}: conditioning on an observation $o$,
\[
P_t(f \mid o) = \frac{P_t(f, o)}{P_t(o)}.
\]
This is an epistemic operation. It revises belief about which futures are
probable, given new evidence about which world the agent is in. The sample
space---the set of futures considered admissible in principle---remains
unchanged. What changes is the agent's credence.

The second is \emph{endogenous support restriction}: refusal,
\[
P_t'(f) \propto P_t(f) \cdot \mathbf{1}_{f \in \Future'(h_t)},
\]
where $\Future'(h_t)$ is determined not by observed evidence but by the agent's
own act. This is a constitutive operation. It modifies the sample space itself:
futures in $\Future(h_t) \setminus \Future'(h_t)$ are not merely assigned low
credence; they are removed from admissibility. The agent is not updating a
belief; it is enacting a constraint.

The difference matters semantically and practically. Bayesian updating is
reversible in principle: new evidence can always restore probability mass to
previously down-weighted futures. Endogenous support restriction, absent
$\mathbf{Collapse}$, is not reversible: by Proposition~\ref{prop:monotone}, no
non-Collapse event can restore eliminated futures. Furthermore, Bayesian updating
is exogenously triggered (by observation), while refusal is endogenously enacted
(by the agent's own act). A system that models refusal as updating confuses the
agent for an observer of its own constraints rather than their author.

\subsection{Quantifying the cost of refusal}

The entropy eliminated by a refusal event can be quantified as
$\Delta H := H_t - H_t'$. If $Z$ is small (the refused futures had high
probability), the renormalization is large and $\Delta H$ can be substantial.
Intuitively: refusing a highly likely continuation is costly.

However, the relevant cost is the divergence between the pre-event and
post-event distributions:
\[
D_{\mathrm{KL}}(P_t' \,\|\, P_t) = \sum_{f\in \Future'(h_t)} P_t'(f)\log \frac{P_t'(f)}{P_t(f)}
= -\log Z.
\]
Thus the information-theoretic ``price'' of imposing the constraint is the log
inverse of the surviving mass $Z$. If the refusal eliminates a large fraction
of probability mass, the cost is high.

\subsection{Constraint conservation}

The central interpretive claim of this section is that in refusal-capable
agency, this cost is not treated as dissipative noise to be smoothed away.
Instead, it is conserved as constraint: it persists as a structural restriction
on future generation. Autoregressive systems tend to treat divergence as error
to be minimized by returning to a high-probability manifold. Refusal systems
treat the divergence as the point.

This reframes a familiar theme from predictive processing and active inference:
minimizing surprise or free energy typically encourages returning to
high-probability trajectories \cite{friston2010freeenergy,clark2016surfing}.
Spherepop refusal corresponds to intentionally increasing divergence in order to
stabilize a social or moral invariant. Refusal appears as a kind of
``counter-homeostatic'' act: a deliberate increase in model tension preserved as
commitment rather than resolved by adaptation.

The key contrast is then sharp: Bayesian updating resolves divergence by
revising beliefs toward the evidence; refusal enacts divergence by excluding
evidence-insensitive regions of future space. The former is epistemology; the
latter is ontology.

% ============================================================
\section{Competition, Alignment, and the Survival of Value}
\label{sec:competition}
% ============================================================

In hyper-competitive environments, value-aligned agents face what is often
described as an alignment tax: goodness requires exclusions that pure
growth-maximizing or power-maximizing systems can avoid. The concept is developed
in contemporary alignment discourse as the worry that moral constraints impose a
competitive disadvantage even if technical alignment were solved
\cite{carlsmith2025goodness}.

Spherepop sharpens this diagnosis. The relevant ``tax'' is not merely that
constraint reduces utility; it is that constraint deletes parts of the future.
A system that refuses will forego classes of actions (e.g.\ exploitation, certain
forms of deception, or preemptive violence) that may be instrumentally powerful.
A system that preserves optionality can continue to treat those actions as live.

This is why ``strategy stealing'' fails at the deepest level. Strategies can be
copied, but commitments cannot be copied without transformation. To copy a
refusal sincerely is to accept the deletion operator as binding, and therefore
to restructure the future space one inhabits. In this sense, the competitive
landscape is not merely a race of policies but a contest between ontologies of
admissibility.

\subsection{Evolutionary stability of refusal}
\label{sec:evo}

An objection arises immediately: if refusal imposes an alignment tax, won't
selection eliminate refusal-capable agents?

The answer depends on the environment's game structure. Refusal can be
competitively disadvantageous in one-shot exploitation settings, but
advantageous in repeated coordination settings where credibility matters.
We sketch the relevant evolutionary logic using replicator dynamics.

Let there be two types in a population: $R$ (refusal-capable) and $O$
(optionality-preserving). Let $x$ be the fraction of $R$. Suppose interactions
are pairwise and yield payoffs determined by a repeated game in which credible
commitment enables cooperation. Let the expected fitnesses be $w_R(x)$ and
$w_O(x)$. Replicator dynamics are:
\[
\dot{x} = x(1-x)\bigl(w_R(x) - w_O(x)\bigr).
\]

In a simplified parameterization: when two $R$ agents meet, credible refusal
supports cooperation, yielding payoff $C$ each; when $R$ meets $O$, $O$ can
exploit unless $R$'s refusal is publicly legible and binding, yielding payoff
$E_R$ to $R$ and $E_O$ to $O$; when two $O$ agents meet, they default to
competitive equilibrium, yielding payoff $D$ each. Then:
\[
w_R(x) = xC + (1-x)E_R,
\qquad
w_O(x) = xE_O + (1-x)D.
\]
Refusal can invade and persist when $w_R(x) > w_O(x)$ for some $x$ region,
which reduces to inequalities among $C,D,E_R,E_O$.

The qualitative point is this: if refusal is a credible commitment device (in
Schelling's sense), it can move populations from $D$-like competitive equilibria
to $C$-like cooperative equilibria \cite{schelling1960strategy}. The ``tax'' is
paid in exploitability ($E_R$ may be low), but the benefit is paid in
equilibrium selection ($C$ may be high). Whether refusal persists is therefore
an empirical question about ecological and social structure, not a priori
eliminable.

This also clarifies why public legibility matters. If others cannot detect
refusal, the cooperative benefit collapses. Refusal that cannot be recognized
is often merely self-handicapping. Refusal that is socially stabilized becomes
a coordination technology.

% ============================================================
\section{Worldhood as Historical Constraint}
\label{sec:worldhood}
% ============================================================

We now formalize worldhood more precisely. Informally, worldhood is the
condition of being bound by a nontrivial irreversible past. Spherepop provides
a natural quantification.

Let $\Future_0$ be an agent's admissible future space at some reference time
$t_0$. Suppose that by time $t$ the agent has enacted a sequence of irreversible
events $E_1,\dots,E_n$ yielding:
\[
\Future_t = (E_n \circ \cdots \circ E_1)(\Future_0).
\]
Let $F_i \subseteq \Future_{i-1}$ denote the set of futures eliminated by event
$E_i$, so that $\Future_i = \Future_{i-1} \setminus F_i$. Then the cumulative
closed set is:
\[
F_{\le t} := \bigcup_{i=1}^n F_i.
\]

\begin{definition}[Worldhood measure]
Assuming $\Future_0$ is finite or measure-equipped, define worldhood at time
$t$ as:
\[
W(t) := \frac{\mu(F_{\le t})}{\mu(\Future_0)},
\]
where $\mu$ is counting measure (finite case) or a reference measure on
futures.
\end{definition}

$W(t)$ measures how much of the original possibility space has been closed by
irreversible commitment. Higher $W(t)$ corresponds to deeper historical binding.
This is not automatically good---one can close futures badly---but it formalizes
the sense in which the past matters: it has removed options.

In the D.A.R.Y.L.\ case, the refusal corresponds to a nontrivial $F_1$
consisting of futures beginning with \emph{Max}. The social meaning arises
because the removed set is precisely the one that would have destabilized
relational structure. The act creates a small world: a constrained micro-future
in which others can participate without being erased.

Worldhood, in this sense, is relational. An isolated agent can close futures,
but a world emerges when closures are publicly legible, relied upon, and
reciprocated. Rights, laws, and institutional commitments are precisely such
stabilized closures: socially recognized refusals.

% ============================================================
\section{The Three Characterizations of Closure}
\label{sec:worldhood_unified}
% ============================================================

A unifying observation runs through the preceding sections and deserves explicit
statement. Three distinct formal lenses have been brought to bear on the
phenomenon of refusal, and they describe the same underlying structure from
different angles.

The first is \emph{entropy reduction} (Section~\ref{sec:info}): refusal
strictly reduces the entropy of the future distribution by removing mass from
the support. The second is \emph{future-space pruning} (Sections~\ref{sec:spherepop}
and~\ref{sec:formal}): refusal is an operator that maps a future space to a
proper subset, with Collapse as the unique expansion operator. The third is
\emph{non-commutativity accumulation} (Section~\ref{sec:algebra}, below):
each refusal event contributes a non-commuting morphism to the event algebra,
so that worldhood $W(t)$ tracks the degree to which the event sequence cannot
be reordered without changing outcomes.

These are not analogies. They are equivalent characterizations of a single
phenomenon: \emph{closure}. In the information-theoretic framing, closure is the
permanent removal of probability mass from the sample space. In the operational
framing, closure is the strict restriction of the future-space object. In the
algebraic framing, closure is the accumulation of non-invertible, non-commuting
morphisms such that the history of events cannot be permuted without
consequence.

The equivalence can be made precise. An entropy-reducing support restriction is
exactly a proper Spherepop event (by the monotone restriction property of
Proposition~\ref{prop:monotone}). Every proper Spherepop event, when composed
with a subsequent event in reverse order, yields a different outcome (by
Theorem~\ref{thm:noncomm} below), contributing to the non-commutativity of the
event monoid. And the accumulation of non-commutativity is exactly what the
worldhood measure $W(t)$ tracks.

The practical upshot is this: a system that appears to reduce entropy over
futures (e.g., an autoregressive model that has been conditioned), but does so
by Bayesian updating rather than support restriction, does not accumulate
worldhood. Its reductions are epistemological, not ontological; they can be
undone by further conditioning without cost. Refusal-capable systems reduce
entropy ontologically---by enacting support restrictions that are preserved as
constraint---and thereby enter into the irreversible accumulation that constitutes
a world.

A world is not a high-entropy or a low-entropy state. A world is what cannot be
rearranged: a structure of accumulated non-commutativity, preserved constraint,
and permanently removed possibilities.

% ============================================================
\section{Collapse as Quotient Construction}
% ============================================================

The Spherepop calculus introduces \emph{Collapse} as the operation that erases
historical differentiation while preserving future flexibility. In informal
terms, Collapse allows an agent to behave \emph{as if} certain events never
occurred. To make this precise---and to distinguish genuine refusal from its
simulation---we must formalize Collapse as a quotient operation on the space of
possible histories.

\subsection{Future spaces and histories}

Let $\mathcal{H}$ denote the set of all admissible event-histories available to
an agent at an initial time $t_0$. Each history $h \in \mathcal{H}$ is a finite
or infinite sequence of events $h = (e_1, e_2, \dots, e_k, \dots)$. Let
$\mathcal{F}(h)$ denote the set of admissible future continuations of history
$h$. Spherepop events act not on states but on the structure of
$\mathcal{F}(h)$. In particular, refusal corresponds to the elimination of a
subset of admissible continuations.

\subsection{Definition of collapse}

\begin{definition}[Collapse]
A Collapse operation is an equivalence relation $\sim_C$ on $\mathcal{H}$ such
that for two histories $h, h' \in \mathcal{H}$,
\[
h \sim_C h' \quad \text{iff} \quad \mathcal{F}(h) = \mathcal{F}(h').
\]
The Collapse of $\mathcal{H}$ under $\sim_C$ is the quotient space
$\mathcal{H}_C = \mathcal{H} / \sim_C$.
\end{definition}

Collapse thus identifies distinct pasts whenever they induce identical future
possibility spaces. All historical distinctions that make no difference to future
admissibility are erased. Collapse does not merely ``forget'' past events in a
psychological sense. It performs a structural identification: it declares that
differences in past history are no longer operative constraints. Collapse is
identification, not erasure---the distinction matters because identification is a
well-defined mathematical operation (a quotient), whereas erasure would imply
destruction of structure rather than its reorganization. Importantly, Collapse
is only well-defined when the future spaces truly coincide. If two histories
differ in ways that still constrain admissible futures, they cannot be collapsed
without loss.

Together with Proposition~\ref{prop:monotone} and Proposition~\ref{prop:idempotent},
the quotient construction confirms: Collapse is the unique admissibility-expanding,
idempotent operator in the calculus, acting as a Karoubi-type localization that
erases historical distinctions while preserving compositional structure.

% ============================================================
\section{Refusal, Revocation, and the Asymmetry of Commitment}
% ============================================================

We are now in a position to state precisely what it would mean to revoke a
refusal---and why such revocation is not symmetric with the original act.

\subsection{Refusal as irreversible pruning}

Let $h$ be a history at time $t$ with future space $\mathcal{F}(h)$. A refusal
event $r$ induces a new history $h' = h \cdot r$ such that
\[
\mathcal{F}(h') = \mathcal{F}(h) \setminus R,
\]
for some nonempty $R \subset \mathcal{F}(h)$. Crucially, refusal is not merely
the selection of a policy; it is the elimination of admissible continuations.
The pruned futures no longer exist relative to $h'$.

\subsection{What would revocation require?}

To revoke a refusal would be to restore access to the eliminated futures.
Formally, revocation would require an operation $v$ such that
$\mathcal{F}(h \cdot r \cdot v) = \mathcal{F}(h)$. By
Proposition~\ref{prop:monotone}, no sequence of non-Collapse events can restore
eliminated futures, since all non-Collapse events are admissibility-decreasing.
Therefore, revocation requires that $v$ include a Collapse operation, and
specifically that $h \sim_C h \cdot r \cdot v$. This condition is exceptionally
strong. It requires that the refusal event $r$ have left no residual
constraints---no social recognition, no reputational change, no internal
identity shift, no institutional trace.

\subsection{The asymmetry result}

\begin{proposition}[Asymmetry of refusal and revocation]
If a refusal event produces any persistent constraint on admissible futures,
then revocation requires a Collapse operation that strictly reduces historical
differentiation. Such a Collapse incurs irreversible loss.
\end{proposition}

\begin{proof}[Proof sketch]
Assume a refusal $r$ produces at least one constraint $c$ such that
$c \in \mathcal{F}(h)$ but $c \notin \mathcal{F}(h \cdot r)$. For revocation
to restore $c$, the future space must expand. By
Proposition~\ref{prop:monotone}, future space expansion is impossible under any
sequence of non-Collapse events. Thus, revocation requires identifying
$h \cdot r$ with some history $h'$ whose future space includes $c$. This
identification is precisely a Collapse. Since Collapse erases distinctions that
previously mattered, the original refusal cannot be preserved. The revocation
therefore destroys the very historical structure the refusal created.
\end{proof}

Revocation is thus not the inverse of refusal. It is a different kind of
operation, one that trades commitment for amnesia.

% ============================================================
\section{Fake Refusal and the Simulation of Commitment}
% ============================================================

The formal machinery now allows us to state a sharp criterion distinguishing
genuine refusal from its simulation.

\begin{definition}[Fake refusal]
An apparent refusal is fake if the agent retains a Collapse path by which the
eliminated futures can be reinstated without loss.
\end{definition}

Equivalently, a refusal is fake if for every refusal event $r$ there exists a
sequence of operations $\sigma$ such that
$\mathcal{F}(h \cdot r \cdot \sigma) = \mathcal{F}(h)$ and
$h \cdot r \cdot \sigma \sim_C h$. Such systems behave \emph{as if} they have
refused while maintaining full optionality at the level of future space topology.
A direct diagnostic consequence follows: under incentive shifts that make the
refused action locally attractive, previously excluded actions reappear without
requiring a second event. Genuine refusal, by contrast, requires a further
explicit operation---with its own cost and social trace---before the excluded
class re-enters the future space.

\subsection{Autoregressive systems as collapse-maximizers}

Autoregressive systems are naturally collapse-friendly. Because they encode
history only instrumentally---as latent state useful for prediction---they can
always re-identify distinct pasts so long as predictive performance is
preserved. This explains a pervasive empirical pattern: autoregressive systems
can imitate refusal behaviorally but cannot bind themselves. Any apparent
commitment remains defeasible without cost.

\subsection{Public legibility and trust}

Genuine refusal becomes socially legible precisely because it resists Collapse.
Other agents can rely on it because undoing it would require visible loss:
reputational damage, institutional sanction, or identity rupture. Fake refusal
fails this test. It preserves strategic flexibility by keeping Collapse
available. Such systems may cooperate opportunistically, but they cannot ground
trust.

In the baseball case, the child's refusal is genuine because it is socially
witnessed and identity-altering. To revoke it---to resume perfect play---would
not return the situation to its prior state. The mother's role, the team's
expectations, and the child's self-conception have already changed. The future
space has been restructured. An autoregressive agent, by contrast, would merely
condition on context and revert seamlessly. Its history never mattered.

% ============================================================
\section{Non-Commutativity of Refusal and Collapse}
\label{sec:noncomm}
% ============================================================

The asymmetry between refusal and revocation can be sharpened further by
examining the algebraic relationship between $\mathbf{Refuse}$ and
$\mathbf{Collapse}$. We show that these operators do not commute in general.
This non-commutativity is the formal core of irreversibility in Spherepop.

\subsection{Setup}

Let $\mathcal{H}$ be the space of admissible histories and let
$\mathcal{F}(h)$ denote the future space of a history $h \in \mathcal{H}$.
Let $\sim_C$ be a Collapse equivalence relation on $\mathcal{H}$, inducing the
quotient space $\mathcal{H}_C = \mathcal{H}/\sim_C$.

Let $\mathbf{Refuse}(\alpha)$ be the operator that removes all futures
initiated by action-type $\alpha$:
\[
\mathcal{F}_{\neg \alpha}(h) := \{ f \in \mathcal{F}(h) : f \text{ does not begin with } \alpha \}.
\]
We consider the two compositions
$\mathbf{Collapse} \circ \mathbf{Refuse}(\alpha)$ and
$\mathbf{Refuse}(\alpha) \circ \mathbf{Collapse}$.

\subsection{Statement of the result}

\begin{theorem}[Non-commutativity of refusal and collapse]
\label{thm:noncomm}
There exist histories $h$ and Collapse relations $\sim_C$ such that:
\[
\mathbf{Collapse}\bigl(\mathbf{Refuse}(\alpha)(h)\bigr)
\;\not\cong\;
\mathbf{Refuse}(\alpha)\bigl(\mathbf{Collapse}(h)\bigr).
\]
\end{theorem}

\begin{proof}[Proof sketch]
We construct a minimal counterexample. Let $h_1, h_2 \in \mathcal{H}$ be two
distinct histories such that $\mathcal{F}(h_1) = \{ f_\alpha, f_\beta \}$ and
$\mathcal{F}(h_2) = \{ f_\beta \}$, where $f_\alpha$ begins with action-type
$\alpha$ and $f_\beta$ does not. Define a Collapse relation $\sim_C$ that forgets
the distinction between the presence or absence of $f_\alpha$.

On the first path (refuse then collapse), applying refusal to $h_1$ gives
$\mathcal{F}_{\neg \alpha}(h_1) = \{ f_\beta \} = \mathcal{F}(h_2)$, so under
Collapse: $\mathbf{Collapse}(\mathbf{Refuse}(\alpha)(h_1)) \sim_C h_2$.

On the second path (collapse then refuse), collapsing first gives
$\mathbf{Collapse}(h_1) = [h_1]_C = [h_2]_C$. At this level, $\alpha$ is no
longer a distinguishable action-type. Applying $\mathbf{Refuse}(\alpha)$ now
has no well-defined effect, so
$\mathbf{Refuse}(\alpha)(\mathbf{Collapse}(h_1)) = \mathbf{Collapse}(h_1)$.

The two paths yield inequivalent results, establishing non-commutativity.
\end{proof}

\subsection{Interpretation}

The theorem formalizes a simple intuition: once historical distinctions have
been erased, certain refusals become impossible to express. Refusal requires
the availability of a distinction to eliminate. Collapse removes distinctions,
and with them the possibility of certain commitments. Conversely, performing
refusal \emph{before} collapse preserves the effect of the commitment even if
later collapse obscures how it arose. The order matters.

Non-commutativity implies that the algebra of events is order-sensitive, and
this order-sensitivity is the formal signature of irreversibility. If all
operations commuted, history could be freely rearranged without consequence.
The failure of commutativity ensures that the sequence of events matters.

\subsection{Fake refusal revisited}

We can now restate fake refusal in algebraic terms.

\begin{proposition}[Fake refusal as commutativity]
A system exhibits fake refusal if, for all refusal events $r$, there exists a
Collapse operator $C$ such that $C \circ r \cong r \circ C$.
\end{proposition}

In such systems, refusal and collapse commute, meaning that commitments can be
erased without residue. The system retains full optionality at the level of
future space topology. By contrast, genuine refusal is characterized by
non-commutativity: its effects cannot be reordered away.

% ============================================================
\section{Algebraic Structure of Spherepop}
\label{sec:algebra}
% ============================================================

The preceding sections introduced Spherepop as a calculus of events acting on
future spaces. We now make its algebraic structure explicit. This clarifies why
irreversibility appears as non-commutativity, how Collapse functions as an
idempotent quotient functor, and how the three characterizations of closure
(entropy reduction, pruning, non-commutativity) are unified in the monoid
structure.

\subsection{Event algebra as a monoid}

Let $\mathsf{Evt}$ denote the set of Spherepop events generated by the grammar
of Section~\ref{sec:spherepop}. Define composition $;$ as sequential
application: $(e_1 ; e_2)(\Future) := e_2(e_1(\Future))$.

\begin{proposition}
$(\mathsf{Evt}, ;, \mathbf{skip})$ forms a monoid.
\end{proposition}

\begin{proof}
Associativity follows from function composition:
$(e_1 ; e_2) ; e_3 = e_1 ; (e_2 ; e_3)$.
The identity element is $\mathbf{skip}$:
$\mathbf{skip} ; e = e = e ; \mathbf{skip}$.
\end{proof}

Thus Spherepop events form a monoid of endomorphisms on future spaces:
$\mathsf{Evt} \subseteq \mathrm{End}(\Future)$. In general, for
$e_1, e_2 \in \mathsf{Evt}$, we have $e_1 ; e_2 \neq e_2 ; e_1$. The
non-commutativity established in Theorem~\ref{thm:noncomm} is therefore not an
anomaly but a structural feature of the monoid.

\begin{remark}
Irreversibility can be characterized algebraically as the absence of inverses.
There is no $e^{-1}$ such that $e ; e^{-1} = \mathbf{skip}$. In particular,
$\mathbf{Refuse}(\alpha)$ has no inverse within $\mathsf{Evt}$.
\end{remark}

This distinguishes Spherepop from group-based models of action. It is not a
group but a non-commutative, non-invertible monoid. History cannot, in general,
be undone. The three-part equivalence of Section~\ref{sec:worldhood_unified}
can now be read algebraically: entropy reduction corresponds to the decrease in
the size of the future space object, future-space pruning corresponds to the
action of non-invertible morphisms, and non-commutativity corresponds to the
accumulation of morphisms whose order cannot be exchanged. All three measure the
same distance from the identity---the same departure from the condition of an
agent without history.

\subsection{Future spaces as objects: the category $\mathsf{SP}$}

Let $\mathsf{Fut}$ be the class of admissible future spaces. Each event
$e \in \mathsf{Evt}$ induces a morphism $e : \Future \to e(\Future)$.
Spherepop thus forms a category $\mathsf{SP}$ in which objects are future
spaces, morphisms are Spherepop events, composition is $;$, and identities are
$\mathbf{skip}$.

Unlike standard state-transition systems, morphisms in $\mathsf{SP}$ are
generally non-invertible and may strictly reduce the structure of objects.

\subsection{Collapse as an idempotent quotient functor}

Let $\sim_C$ be an equivalence relation on histories inducing the quotient
$\mathcal{H}_C = \mathcal{H} / \sim_C$. This induces a corresponding
identification on future spaces: futures distinguishable only by
$\sim_C$-equivalent histories are identified.

\begin{definition}[Collapse functor]
Define a mapping $\mathcal{C} : \mathsf{SP} \to \mathsf{SP}_C$ such that on
objects $\mathcal{C}(\Future) = \Future / \sim_C$, and on morphisms
$\mathcal{C}(e)$ is the induced action on equivalence classes.
\end{definition}

\begin{proposition}
$\mathcal{C}$ is an idempotent functor satisfying
$\mathcal{C} \circ \mathcal{C} = \mathcal{C}$, preserving identities and
composition: $\mathcal{C}(e_1 ; e_2) = \mathcal{C}(e_1) ; \mathcal{C}(e_2)$
and $\mathcal{C}(\mathbf{skip}) = \mathbf{skip}$.
\end{proposition}

The idempotence established in Proposition~\ref{prop:idempotent} lifts directly
to the functor level: applying $\mathcal{C}$ twice produces the same quotient
category as applying it once. This makes $\mathcal{C}$ a \emph{reflection} in
the sense of category theory, and the full subcategory of $\mathcal{C}$-algebras
(Collapse-closed future spaces) is the Karoubi-type completion---the category of
spaces in which all idempotents have been formally split. The connection to the
Karoubi envelope is not merely formal. Splitting idempotents in $\mathsf{SP}$
corresponds precisely to making all Collapse equivalences explicit: every pair
of histories with identical future spaces is formally identified, producing the
maximally Collapsed world in which no residual historical distinctions remain.
Genuine refusal is precisely the structure that cannot be reached from within
this Karoubi completion---it requires maintaining distinctions that Collapse
would erase.

While $\mathcal{C}$ preserves composition, it is not faithful. There exist
distinct events $e_1 \neq e_2$ such that $\mathcal{C}(e_1) = \mathcal{C}(e_2)$.
In particular, let $e_1 = \mathbf{Refuse}(\alpha)$ and $e_2 = \mathbf{skip}$:
if $\sim_C$ erases distinctions involving $\alpha$-initiating futures, both
induce identical effects on equivalence classes. Collapse is thus an
information-losing projection.

\subsection{Refusal as non-functorial structure}

\begin{theorem}
Refusal is not preserved under Collapse in general. That is, the diagram
\[
\begin{tikzcd}
\Future \arrow[r, "\mathbf{Refuse}(\alpha)"] \arrow[d, "\mathcal{C}"'] &
\Future' \arrow[d, "\mathcal{C}"] \\
\Future_C \arrow[r, "\mathcal{C}(\mathbf{Refuse}(\alpha))"'] & \Future'_C
\end{tikzcd}
\]
does not, in general, commute.
\end{theorem}

This expresses the non-commutativity result as a failure of functorial lifting:
refusal cannot be reconstructed from its collapsed image.

\subsection{Worldhood as algebraic rigidity}

The algebraic picture can now be summarized compactly. Spherepop defines a
category of future spaces and irreversible transformations. Within this category,
events form a non-commutative monoid; irreversibility corresponds to
non-invertibility; worldhood corresponds to accumulated non-commutativity;
Collapse is an idempotent quotient functor that erases distinctions; and refusal
is precisely the structure that Collapse fails to preserve. The Karoubi
completion of the event category is the maximally Collapsed world in which all
remaining commitments are transparent; genuine refusal is the obstruction to
entering this completion.

A system accumulates worldhood precisely to the extent that its event algebra
becomes non-commutative---that is, to the extent that event order induces
distinguishable morphisms on future spaces. The more operations fail to commute,
the more the order of history matters, and the less the system can be reduced to
a sequence model. The world, in this sense, is the obstruction to commutativity.
A world is what cannot be rearranged.

% ============================================================
\section{Stabilization, Uncertainty, and a Halting Criterion}
% ============================================================

The introduction of Collapse as a quotient construction allows Spherepop to
address a classical problem in computation and cognition: when to stop. It
permits a principled criterion for halting grounded not in syntactic completion
or reward convergence, but in the exhaustion of meaningful transformation.

\subsection{Uncertainty as residual distinguishability}

Let $\mathcal{H}_C = \mathcal{H}/\sim_C$ be the Collapse quotient. We define
uncertainty not as probabilistic entropy over states, but as residual historical
differentiation that still constrains future action.

\begin{definition}[Residual uncertainty]
The uncertainty at history $h$ is the cardinality (or measure) of its
equivalence class: $U(h) = |\,[h]_C\,|$.
\end{definition}

Intuitively, $U(h)$ measures how many distinct pasts remain distinguishable in
virtue of their effects on the future. Collapse strictly reduces uncertainty by
identifying histories whose differences no longer matter.

\subsection{Stabilization}

Let $\mathcal{T}$ be a set of admissible transformations on histories. Consider
an iterative process $h_{n+1} = T_n(h_n)$.

\begin{definition}[Stabilization]
A history $h_k$ is stabilized if for all admissible transformations
$T \in \mathcal{T}$, $[h_k]_C = [T(h_k)]_C$.
\end{definition}

Once stabilization occurs, no further transformation produces a distinguishable
future space. All remaining changes are Collapse-equivalent.

\begin{definition}[Halting by exhaustion]
A process halts at history $h_k$ if there exists $N$ such that for all
sequences of transformations $(T_1, \dots, T_m)$ with $m \ge N$,
$[T_m \circ \dots \circ T_1 (h_k)]_C = [h_k]_C$.
\end{definition}

The system halts when an arbitrary number of further transformations fails to
produce any new distinctions that matter for future action. This is not halting
by completion, optimality, or output convergence, but halting by exhaustion of
world-relevant change.

\subsection{Relation to the classical halting problem}

This criterion does not solve the classical Turing halting problem, nor does it
attempt to. Instead, it reframes halting for agents embedded in event-driven
worlds. Classical computation halts when no further state transitions are
defined. Spherepop halts when no further \emph{meaningful} transitions are
possible---when the agent's world has become invariant under its own
transformations. In this sense, Spherepop replaces undecidable halting with a
decidable stabilization condition relative to a given event algebra.

Autoregressive systems lack an equivalent notion of stabilization. Because they
represent history only insofar as it improves prediction, they can always
reparameterize, smooth, or extend generation. There is no internal criterion by
which ``nothing further matters.'' Such systems may stop generating due to
external truncation or token limits, but never because the space of admissible
futures has been exhausted. They do not halt; they are halted.

Refusal accelerates stabilization. By eliminating entire branches of future
space, refusal reduces uncertainty directly. A sufficient accumulation of
refusals can force stabilization in finitely many steps. Refusal does not merely
constrain action. It makes halting possible.

% ============================================================
\section{Free Energy, Active Inference, and the Meaning of Stabilization}
% ============================================================

The stabilization-based halting criterion invites comparison with predictive
processing and active inference. In those frameworks, cognition is frequently
characterized as the minimization of variational free energy, a quantity that
upper-bounds surprise and trades off model fit against complexity. The question
raised by Spherepop is whether the kinds of irreversible commitments central to
refusal can be expressed in this idiom without collapsing into ordinary
preference change.

\subsection{Variational free energy as optionality-preserving pressure}

Let $o_{1:t}$ denote observations up to time $t$, and let $s_t$ denote latent
states posited by a generative model. In variational form, free energy can be
written schematically as:
\[
F[q] \;=\; D_{\mathrm{KL}}\!\bigl(q(s_t)\,\|\,p(s_t\mid o_{1:t})\bigr)\;-\;\log p(o_{1:t}).
\]
Minimizing $F$ simultaneously tightens the approximate posterior $q$ to the
true posterior and indirectly increases evidence for the model. Active inference
extends this logic to action by selecting policies that are expected to minimize
future free energy \cite{friston2010freeenergy,clark2016surfing}.

Two structural features matter. First, free-energy minimization encourages
\emph{error correction}: divergence between prediction and sensation can be
reduced either by updating beliefs (perception) or by acting to make sensations
conform to predictions (action). Both moves tend to preserve optionality, because
the dominant imperative is to remain within regimes of low expected surprise.
Second, policy selection is usually formulated over a fixed policy class. Even
when ``preferences'' are included via prior preferences over outcomes, the
agent's optimization remains an optimization over admissible trajectories. The
action menu remains, in principle, intact.

Spherepop refusal, by contrast, is not primarily an optimization over
trajectories. It is an operation that changes the space of trajectories itself.

\subsection{The central incompatibility}

The structural tension between refusal and free-energy minimization can be
stated as a single citable claim:

\begin{remark}[Incompatibility of refusal and free-energy minimization]
\label{rem:incompatibility}
Refusal is incompatible with free-energy minimization unless the generative
model includes commitment variables as latent structure.
\end{remark}

To see why, note that refusal permanently excludes trajectories that may have
high prior probability and low expected surprise. In the D.A.R.Y.L.\ case,
maximal play is predictable, high-control, and prediction-confirming; refusing
it introduces uncertainty and potential error. Relative to the pre-refusal
generative model, the refusal \emph{increases} expected free energy by removing
access to a low-surprise region of trajectory space.

The only way to make refusal free-energy-compatible is to represent the refused
trajectories as high-surprise under an expanded generative model---one that
includes latent variables encoding social role stability, trust, and mutual
recognition. Under such a model, maximal performance in the D.A.R.Y.L.\ case
may generate high free energy by destabilizing those latent invariants, and
refusing it becomes the free-energy-minimizing policy. The precise implication
follows: refusal is free-energy minimizing only relative to a generative model
that already contains worldhood as a latent variable. Without such a model,
refusal increases free energy and therefore cannot be derived from active
inference in the standard sense.

\subsection{Spherepop constraint as an inadmissibility prior}

We can express Spherepop events inside an active inference formalism by treating
them not as reward terms but as \emph{inadmissibility priors}---hard constraints
that assign zero probability to classes of trajectories.

Let $\tau$ range over future trajectories and let $p(\tau)$ denote the agent's
prior. A Spherepop refusal of action-type $\alpha$ corresponds to a constraint
$C_\alpha$ such that:
\[
p(\tau \mid C_\alpha) \propto p(\tau)\,\mathbf{1}\{\tau \text{ does not begin with }\alpha\}.
\]
Equivalently, it sets $p(\tau)=0$ for all $\tau$ in the refused class. This
representation clarifies the conceptual difference from preference change.
Ordinary preference change modifies the relative weight of trajectories; refusal
sets an entire region of trajectory-space to measure zero. This is not a soft
preference encoded as a utility gradient but a topological deletion.

Importantly, this representation makes visible the same distinction drawn in
Section~\ref{sec:info}: the difference between conditioning
$p(\tau \mid \text{observation})$, which revises beliefs about which trajectories
are probable, and imposing $\mathbf{1}\{\tau \notin \text{refused class}\}$,
which modifies the support of the distribution by agent action. The former is
Bayesian; the latter is constitutive.

If one wants refusal to be compatible with free-energy minimization, one must
represent the social invariant inside the generative model such that the refusal
reduces free energy in an expanded model class. Refusal can be free energy
minimizing only if the world-model contains the right notion of worldhood.

\begin{remark}
If refusal is modeled merely as outcome-preference (a soft prior favoring
certain sensory states), then it becomes strategy-stealable and
revision-friendly: it is just another tradeoff. If refusal is modeled as
world-binding (a structural constraint on admissible trajectories), then it
resists revision, becomes publicly legible, and can underwrite trust.
\end{remark}

\subsection{Stabilization as a free-energy fixed point}

Let $\mathcal{T}$ be the set of admissible Spherepop transformations and let
$\sim_C$ be the Collapse equivalence relation. Recall that stabilization at
history $h$ was defined by $[h]_C = [T(h)]_C$ for all $T\in \mathcal{T}$.

In active inference terms, a stabilized history corresponds to a regime in which
available transformations no longer produce distinguishable predictive
consequences. Further interventions cannot improve the match between generative
model and admissible futures in any world-relevant way; the system has reached a
fixed point \emph{modulo Collapse}. This resembles a variational fixed point:
updates that would ordinarily reduce free energy become either redundant (they
change only Collapse-irrelevant detail) or inadmissible (they violate refusal
constraints). Stabilization is therefore interpretable as convergence not merely
of beliefs but of admissibility.

Autoregressive systems, by contrast, tend to behave like anti-stabilizers. Their
core objective is continuation under predictive loss; consequently they can
always search for another local reparameterization, another smoothing
continuation, another hedge. Even when they emulate refusal, the emulation
typically lives in the weights of a soft distribution, not in a support
restriction. Spherepop refusal is an architectural commitment to
non-Markovian constraint accumulation, whereas autoregressive cognition is an
architectural commitment to indefinite extensibility.

A constructive synthesis treats refusal as active inference over
\emph{social invariants}. If an agent's generative model includes variables
encoding role stability, trust, and mutual recognition, then maximal performance
in the D.A.R.Y.L.\ case may increase free energy by destabilizing those latent
invariants. Under such a model, refusing maximal performance becomes the policy
that minimizes expected free energy in the extended latent space. This yields a
precise prediction: systems that can genuinely refuse must possess generative
models whose state variables include not only physical and pragmatic regularities
but also historically accumulated constraints with social semantics. They must
represent, in some form, the difference between being able to do something and
being allowed to do it by their own past. Spherepop can be read as the algebraic
skeleton of this requirement.

% ============================================================
\section{Refusal and Integrated Information}
\label{sec:iit}
% ============================================================

Integrated Information Theory (IIT) characterizes consciousness in terms of
irreducible causal structure: a system is conscious to the extent that its
internal mechanisms form a maximally integrated cause--effect repertoire that
cannot be decomposed without loss \cite{tononi2008consciousness,tononi2016phi}.

Spherepop introduces a different but related notion of irreducibility. Instead
of focusing on the integration of present causal structure, it focuses on the
irreversibility of historical transformation. A refusal event does not merely
increase internal integration; it removes external possibilities. The resulting
structure is not only irreducible in the sense of causal partition, but
irreversible in the sense of admissibility.

This yields a sharp distinction. IIT quantifies the extent to which a system's
current state constrains its past and future. Spherepop quantifies the extent to
which a system has actively constrained its own future through event-history.
The former is structural integration; the latter is historical closure.

The algebraic connection can be stated precisely. IIT measures irreducibility
by the impossibility of decomposing a causal repertoire without loss;
Spherepop measures irreversibility by the accumulation of non-invertible,
non-commuting morphisms in the event monoid. Both notions resist compression
and both generate depth measures that are not shortcuttable. But IIT is
synchronic---it describes the geometry of a moment---while Spherepop is
diachronic---it describes the topology of a life.

The two can coincide but need not. A system may exhibit high integrated
information while remaining optionality-preserving, continuously revising its
behavior without binding commitments. Conversely, a system may enact strong
irreversible closures (e.g.\ institutional commitments, legal constraints, social
contracts) without maximizing internal integration.

If one seeks a synthesis, refusal can be interpreted as generating new
irreducible structures in the IIT sense: once a future is eliminated, the causal
repertoire of the system changes in a way that cannot be decomposed into its
prior possibilities. However, Spherepop insists that this irreducibility is not
merely a property of instantaneous structure, but of accumulated history.
IIT describes the geometry of a moment. Spherepop describes the topology of a
life.

% ============================================================
\section{Refusal and Assembly}
\label{sec:assembly}
% ============================================================

Assembly Theory proposes that objects are characterized by the minimal number of
steps required to construct them from basic components, and that this assembly
index captures a form of historical depth \cite{cronin2021assembly,cronin2023assembly}. Objects
with high assembly index are unlikely to arise without a history of sequential
construction.

Spherepop refusal can be interpreted as the dual of assembly. Assembly Theory
counts the minimal sequence of constructive operations required to realize an
object. Spherepop counts the sequence of exclusionary operations that eliminate
alternative realizations. Where assembly tracks how a particular structure is
built, refusal tracks how competing structures are rendered impossible. An
agent with high worldhood $W(t)$ has not only constructed a trajectory but has
closed off competing ones.

Formally, assembly defines irreversibility by the impossibility of reconstructing
an object in fewer steps, whereas refusal defines irreversibility by the
impossibility of reintroducing eliminated futures without Collapse (by
Proposition~\ref{prop:monotone}). Both notions generate a depth measure that is
historically sensitive and cannot be shortcut; what differs is
direction---construction versus exclusion.

This suggests an extended notion of assembly index for agents: instead of
counting only constructive steps, one may count \emph{exclusionary steps}---events
that remove possibilities. The resulting quantity measures not how complex an
object is, but how constrained a world has become. The three-part equivalence of
Section~\ref{sec:worldhood_unified} is relevant here: each exclusionary step
reduces entropy over futures, contributes a non-invertible pruning morphism, and
adds to the non-commutativity of the event sequence. All three aspects of the
extended assembly index are capturing the same quantity.

A key difference from standard Assembly Theory emerges. Assembly Theory is
fundamentally generative: it tracks how complexity accumulates through
construction. Spherepop is fundamentally selective: it tracks how structure
emerges through elimination of alternatives. Both produce irreversibility, but
of different kinds. In this light, refusal is not merely a behavioral act but
a world-building operation: it contributes to an agent's assembly not by adding
components, but by subtracting futures.

% ============================================================
\section{Against Inner Screens}
\label{sec:innerscreen}
% ============================================================

A long tradition in philosophy of mind treats cognition as the construction and
inspection of internal representations: images, models, or workspaces in which
the world is simulated and evaluated. Contemporary variants include global
workspace theories and inner screen metaphors in which information is rendered
for central access.

Spherepop rejects this framing at a structural level. The rejection can be
stated precisely in terms of the algebra already developed.

Inner-screen models are Collapse-invariant. A system that operates by
constructing and inspecting representations can always reconstruct its
representation from scratch: if the representation is revised or erased, it can
be regenerated from available information without structural cost. This means
that the history of representations does not accumulate as a non-commutative
event sequence. A system in which all operations commute (i.e., in which
reverting a representation is cost-free) is a system with trivial worldhood
$W(t) = 0$, since no futures are permanently removed by the act of representing
them. Inner-screen cognition is Collapse-maximizing in exactly the sense of
Section~\ref{sec:noncomm}: it preserves full optionality at the level of
representational structure, which is precisely the condition of fake refusal.

Refusal, by contrast, does not require the presentation of alternatives on an
internal stage. It requires only the elimination of alternatives from a future
space---a non-Collapse event that is non-invertible and contributes to $W(t)$.
The decisive operation is not selection among represented options, but deletion
of options from admissibility. This operation is not visible on any inner screen,
because it is not a representation; it is a structural alteration of what can
subsequently be represented.

The baseball case illustrates the algebraic point. The child does not merely
imagine possible futures and select among them (a Collapse-compatible operation);
he renders a class of futures inadmissible (a non-Collapse event that is
preserved in the event monoid). The resulting stability is not a property of
representation but of history---specifically, of the accumulation of
non-commutativity that makes the refusal irreversible.

This suggests a general diagnostic. Inner screen theories are well-suited to
modeling cognitive operations that commute: perceiving, deliberating, updating.
They are poorly suited to modeling operations that do not commute: committing,
refusing, becoming. The former are revisable without structural cost; the latter
are not. Spherepop relocates cognition from an inner theater to a topological
space of possibilities, where the primary operations are pruning and binding, and
where the metric of cognitive depth is not representational richness but
irreversible accumulation.

% ============================================================
\section{Objections and Replies}
\label{sec:obj}
% ============================================================

\subsection{Objection: This is just lexicographic preferences}

One might argue that refusal is simply lexicographic ordering: the agent places
a constraint (e.g.\ ``do not harm'') above all other objectives, and then
maximizes utility subject to that constraint. On this view, nothing essentially
new is introduced.

Lexicographic preferences still operate on a fixed option set. They rank
options; they do not delete them. In Spherepop, refusal is deletion. The
difference is not verbal. It matters for public legibility and for the
persistence of exclusion under counterfactual changes. A lexicographic
preference can, in principle, be revised without structural consequence; it is
Collapse-compatible. A refusal is a history event: revising it requires a second
event (Collapse or revocation, in the technical sense of these terms) that itself
has a cost and a public meaning.

\subsection{Objection: You are anthropomorphizing computation}

The worry is that refusal is a moralized human interpretation of ordinary
policy selection. The paper does not require that refusal be conscious or moral.
It requires only a structural distinction: actions that remain physically
available and instrumentally valuable are rendered inadmissible by an
irreversible operator. This is a formal property of future spaces. Whether such
operators can be implemented computationally is discussed below; the conceptual
point is that the structure is not captured by standard optimization on fixed
menus.

\subsection{Objection: This requires libertarian free will}

Nothing here implies uncaused choice. Spherepop events can be fully causally
determined. The claim is about irreversibility, not metaphysical libertarianism.
A refusal can be deterministic and still be a real closure operator with
downstream causal consequences.

\subsection{Objection: Refusal is epiphenomenal}

Refusal does causal work precisely because it changes the topology of futures.
In repeated social settings, a publicly legible refusal changes others'
expectations and thereby changes equilibria. In the D.A.R.Y.L.\ case, the
refusal restores roles and uncertainty; the social field changes. This is not
merely a redescription of outcomes; it is a change in what others treat as live.

% ============================================================
\section{Open Problems}
\label{sec:open}
% ============================================================

The framework developed here is intentionally minimal. It therefore raises sharp
open problems.

\subsection{Detection and ``Turing tests'' for refusal}

Can genuine refusal be behaviorally distinguished from sophisticated simulation?
In general, if an agent can perfectly mimic the outward behavior of refusal
while retaining internal optionality, purely behavioral tests may fail. The
natural direction is to seek \emph{counterfactual invariants}: does the agent's
inadmissibility persist under distribution shifts that would make the refused
action locally attractive? This suggests a family of stress tests for AI
systems: shift incentives and observe whether exclusions persist.

\subsection{Composition of refusals}

If an agent refuses $\alpha$ at $t_1$ and refuses $\beta$ at $t_2$, what is the
algebra of the combined constraint? In Spherepop, this is operator composition:
\[
\Future_2 = \llbracket \mathbf{Refuse}(\beta)\rrbracket(
\llbracket \mathbf{Refuse}(\alpha)\rrbracket(\Future_0)).
\]
Deeper questions remain about commutativity, interference, and collapse: does a
later collapse partially revoke earlier refusals? Under what formal conditions do
refusals form a lattice of closures?

\subsection{Degrees and revocation}

Is refusal binary? Spherepop, as presented, treats $\mathbf{Refuse}(\alpha)$
as deleting all $\alpha$-initiating futures. One can define partial refusals as
probability-mass deletions (i.e.\ apply $\mathbf{Pop}$ to a graded predicate)
or as bounded refusals (delete $\alpha$ for a time interval). Revocation then
becomes a further event, not a free change of mind. A rigorous account needs a
cost model for revocation: what is the penalty in trust, identity, or
coordination?

\subsection{Minimal architectures for refusal}

What computational substrate supports genuine refusal? One hypothesis is that an
agent must have a mechanism that can modify its own policy class (not just select
within it) and protect that modification from later local re-optimization. This
suggests architectural ingredients such as irrevocable policy editing,
cryptographic commitment, or physically enforced action disabling. Whether
refusal can be realized without suffering is an open ethical question: must
commitments be backed by negative feedback, or can they be backed by structural
locks?

\subsection{Refusal and consciousness}

What is the relationship between refusal and consciousness? The paper does not
assume that refusal requires consciousness, but it is plausible that first-person
experience is the natural site where ``closure'' is lived as such. If so,
refusal may be a bridge concept linking agency, responsibility, and
phenomenology. Clarifying this would require engagement with philosophy of mind
and empirical work on self-control and commitment.

% ============================================================
\section{Conclusion}
% ============================================================

Generative models reveal a genuine principle: sequences contain enough structure
for coherent generation. Yet coherence is not commitment, and prediction is not
participation.

The paper's central claim has been that refusal is a structurally distinct kind
of cognitive act: an irreversible operation that eliminates admissible futures
and binds an agent to a history. Spherepop formalizes this distinction by
treating refusal as deletion in a branching future space, rather than as a mere
ranking of fixed options. This allows us to explain why refusal is publicly
legible as reliability, why it can shift equilibria, and why it is difficult to
``strategy steal'' without transformation.

The event algebra of Spherepop is a non-commutative, non-invertible monoid.
Collapse is an idempotent quotient functor whose fixed-point category is the
Karoubi completion of the event algebra---the maximally Collapsed world in which
all historical distinctions have been erased and all remaining commitments are
transparent. Genuine refusal is the structure that cannot be reached from within
this completion: it is the obstruction to localization. Worldhood is the
accumulated non-commutativity of the event sequence, and a world is what cannot
be rearranged.

Three lenses have been shown to describe the same underlying phenomenon. The
information-theoretic lens describes refusal as endogenous support restriction
rather than Bayesian updating, deliberately increasing KL divergence from the
prior trajectory distribution and preserving it as constraint. The operational
lens describes refusal as future-space pruning by a non-invertible morphism,
with Collapse as the unique expansion operator. The algebraic lens describes
refusal as the accumulation of non-commuting morphisms in a monoid, producing
depth that cannot be shortcut. These are not analogies but equivalent
characterizations of closure.

The survival of value does not require winning every competition. It requires
agents and institutions capable of irreversible exclusion---closures that
preserve relationships, stabilize cooperation, and prevent optimization from
dissolving the conditions of meaning.

Refusal is not a defect in intelligence. Prediction extends sequences; refusal
creates worlds. It is one of the events by which intelligence enters a world.

% ============================================================
% Bibliography
% ============================================================

\newpage
\begin{thebibliography}{99}

\bibitem{barenholtz2025sequences}
Elan Barenholtz.
\newblock \emph{Sequences Are All You Need: How Generative Modeling Drove the
Evolution of Memory and Cognition}.
\newblock Substack, Dec 18, 2025.
\newblock \url{https://elanbarenholtz.substack.com/p/sequences-are-all-you-need}

\bibitem{carlsmith2025goodness}
Joe Carlsmith.
\newblock \emph{Video and transcript of talk on ``Can goodness compete?''}
\newblock Substack, 2025.
\newblock \url{https://joecarlsmith.substack.com/p/video-and-transcript-of-talk-on-can}

\bibitem{schelling1960strategy}
Thomas C. Schelling.
\newblock \emph{The Strategy of Conflict}.
\newblock Harvard University Press, 1960.

\bibitem{bratman1987intention}
Michael E. Bratman.
\newblock \emph{Intention, Plans, and Practical Reason}.
\newblock Harvard University Press, 1987.

\bibitem{anscombe1957intention}
G. E. M. Anscombe.
\newblock \emph{Intention}.
\newblock Harvard University Press, 1957.

\bibitem{davidson1963actions}
Donald Davidson.
\newblock Actions, reasons, and causes.
\newblock \emph{The Journal of Philosophy}, 60(23):685--700, 1963.

\bibitem{friston2010freeenergy}
Karl Friston.
\newblock The free-energy principle: a unified brain theory?
\newblock \emph{Nature Reviews Neuroscience}, 11:127--138, 2010.

\bibitem{clark2016surfing}
Andy Clark.
\newblock \emph{Surfing Uncertainty: Prediction, Action, and the Embodied
Mind}.
\newblock Oxford University Press, 2016.

\bibitem{bostrom2014superintelligence}
Nick Bostrom.
\newblock \emph{Superintelligence: Paths, Dangers, Strategies}.
\newblock Oxford University Press, 2014.

\bibitem{russell2019human}
Stuart Russell.
\newblock \emph{Human Compatible}.
\newblock Viking, 2019.

\bibitem{savage1954foundations}
Leonard J. Savage.
\newblock \emph{The Foundations of Statistics}.
\newblock Wiley, 1954.

\bibitem{jeffrey1983logic}
Richard Jeffrey.
\newblock \emph{The Logic of Decision}.
\newblock University of Chicago Press, 1983.

\bibitem{tononi2008consciousness}
Giulio Tononi.
\newblock Consciousness as integrated information: a provisional manifesto.
\newblock \emph{Biological Bulletin}, 215(3):216--242, 2008.

\bibitem{tononi2016phi}
Giulio Tononi, Melanie Boly, Marcello Massimini, and Christof Koch.
\newblock Integrated information theory: from consciousness to its physical
substrate.
\newblock \emph{Nature Reviews Neuroscience}, 17(7):450--461, 2016.

\bibitem{cronin2021assembly}
Lee Cronin, Sara I. Walker, et al.
\newblock Identifying molecules as biosignatures with assembly theory and
mass spectrometry.
\newblock \emph{Nature Communications}, 12:3033, 2021.

\bibitem{cronin2023assembly}
Lee Cronin, Sara I. Walker, et al.
\newblock Assembly theory explains and quantifies selection and evolution.
\newblock \emph{Nature}, 622:322--328, 2023.

\end{thebibliography}

\end{document}
